% engines: lualatex
% expect-warnings: 0
% expect-errors: 0
\documentclass{article}
\usepackage[margin=1cm]{geometry}
\usepackage[kernel = lua]{../../latex/pgf-go}
\newcommand\case[1]{\par\medskip\noindent\textbf{#1}\par\nopagebreak}
\pagestyle{empty}

% Draws frames of \sgfanimation one by one (as in sgf-animation.tex).
\ExplSyntaxOn
\NewDocumentCommand \frames {mmm} {
	\group_begin:
		\pgfgo_sgfanimation_setup:nn {#1} {#2}
		\clist_map_inline:nn {#3} {\pgfgo_sgfanimation_frame:n {##1}\quad}
	\group_end:
}
\ExplSyntaxOff
\begin{document}

% far-branch.sgf: 19x19. Main line Q16, R14, O17 (upper right); the commented variation 2 plays
% D4, F3, C6 (lower left). partial = auto uses only the lines that are drawn. The label a at C17 (move
% 2, as in Kogo's Joseki Dictionary) is not drawn, so it does not count.

\case{sgf diagrams, main line only: the upper right}
\sgfdiagrams[partial = auto, goban = {scale = 0.4}]{far-branch.sgf}

\case{sgf diagrams, to = 2: the label a at C17 is drawn (the whole width); with marks = false, it is not}
\sgfdiagrams[partial = auto, to = 2, goban = {scale = 0.4}]{far-branch.sgf}
\sgfdiagrams[partial = auto, to = 2, marks = false, goban = {scale = 0.4}]{far-branch.sgf}

\case{sgf diagrams, variation = 2 and variation = \{D4\}: the lower left}
\sgfdiagrams[partial = auto, variation = 2, goban = {scale = 0.4}]{far-branch.sgf}
\sgfdiagrams[partial = auto, variation = {D4}, goban = {scale = 0.4}]{far-branch.sgf}

\case{sgf diagrams with variations: main line and variation, the whole board in both}
\sgfdiagrams[partial = auto, variations, goban = {scale = 0.3}]{far-branch.sgf}

\case{sgf animation, variation = 2: frames 0 and 3 in the lower left}
\frames{partial = auto, variation = 2, goban = {scale = 0.4}}{far-branch.sgf}{0, 3}

\end{document}
